% =====================================================================
%  CARNET D'ENTRAÎNEMENT — FICHE 8 — THÈME 1 — NIVEAU MOYEN — CORRIGÉ
% =====================================================================
\documentclass[11pt,a4paper]{article}
\usepackage[utf8]{inputenc}
\usepackage[T1]{fontenc}
\usepackage{lmodern}
\usepackage[french]{babel}
\usepackage{amsmath,amssymb,mathtools}
\usepackage{icomma}
\usepackage{siunitx}
\sisetup{output-decimal-marker={,},per-mode=power,group-separator={\,},group-minimum-digits=5}
\usepackage[version=4]{mhchem}
\usepackage{xcolor}
\usepackage{tikz}
\usepackage[most]{tcolorbox}
\usepackage{enumitem}
\usepackage{array}
\usepackage{fancyhdr}
\usepackage[hidelinks]{hyperref}
\usetikzlibrary{arrows.meta,positioning,calc}
\usepackage[a4paper,margin=2.1cm,headheight=26pt,headsep=10pt]{geometry}
\usepackage{titlesec}

\definecolor{primary}{RGB}{150,98,162}
\definecolor{primarydark}{RGB}{92,52,110}
\definecolor{excol}{RGB}{0,135,90}

\tcbset{boxbase/.style={enhanced,breakable,boxrule=0.9pt,arc=2.5pt,left=3mm,right=3mm,top=2.3mm,bottom=2.3mm,fonttitle=\bfseries,coltitle=white,toptitle=0.6mm,bottomtitle=0.6mm}}
\newtcolorbox[auto counter]{correction}[1][]{boxbase,colback=excol!4,colframe=excol,title={\textsc{Corrigé}~\thetcbcounter\ \ifx\relax#1\relax\relax\else\textmd{--- \itshape #1}\fi}}

\newcommand{\dH}{\Delta H}
\newcommand{\dU}{\Delta U}
\newcommand{\rep}[1]{\textcolor{primarydark}{\textbf{#1}}}

\pagestyle{fancy}
\fancyhf{}
\fancyhead[L]{\small\textcolor{primarydark}{\textbf{Fiche 8} — Thème 1 : Premier principe et enthalpie}}
\fancyhead[R]{\small\textcolor{primarydark}{\textsc{Moyen — corrigé}}}
\fancyfoot[C]{\small\thepage}
\renewcommand{\headrulewidth}{0.8pt}
\renewcommand{\headrule}{\hbox to\headwidth{\color{primary}\leaders\hrule height \headrulewidth\hfill}}
\setlength{\parindent}{0pt}\setlength{\parskip}{3pt}

\begin{document}

\begin{center}
\begin{tikzpicture}
\node[rectangle,rounded corners=7pt,fill=excol,text=white,minimum width=16cm,
      minimum height=1.1cm,align=center,inner sep=6pt]
  {{\large\bfseries Thème 1 — Premier principe et enthalpie}\\[1pt]{\small Niveau \textsc{Moyen} — corrigé détaillé}};
\end{tikzpicture}
\end{center}

\begin{correction}[travail des forces de pression]
On convertit les volumes en \si{\cubic\metre} puis on applique $W=-p\,\Delta V$.
\begin{enumerate}[label=\alph*),leftmargin=2em,topsep=2pt,itemsep=3pt]
  \item $\Delta V = (5{,}0-2{,}0)=\SI{3.0}{\litre}=\SI{3.0e-3}{\cubic\metre}$.
        \[ W = -p\,\Delta V = -(\SI{1.0e5}{\pascal})(\SI{3.0e-3}{\cubic\metre}) = \rep{\SI{-300}{\joule}}. \]
        $W<0$ : en se dilatant, le gaz \rep{cède} du travail (il pousse l'atmosphère).
  \item $\Delta V = (1{,}0-4{,}0)=\SI{-3.0}{\litre}=\SI{-3.0e-3}{\cubic\metre}$.
        \[ W = -(\SI{2.0e5}{\pascal})(\SI{-3.0e-3}{\cubic\metre}) = \rep{\SI{+600}{\joule}}. \]
        $W>0$ : comprimé, le gaz \rep{reçoit} du travail de l'extérieur.
\end{enumerate}
\end{correction}

\begin{correction}[relier $\dH$, $W$ et $\dU$]
\begin{enumerate}[label=\alph*),leftmargin=2em,topsep=2pt,itemsep=2pt]
  \item La réaction libère de la chaleur : $\dH = \rep{\SI{-572}{\kilo\joule}}$, elle est \rep{exothermique}.
  \item $\dU = \dH + W = -572 + 8 = \rep{\SI{-564}{\kilo\joule}}$.
  \item Les réactifs comptent $3\ \si{mol}$ de gaz ($2\,\ce{H2}+\ce{O2}$), les produits $0$ ($\ce{H2O}$
        liquide) : le volume \emph{diminue} fortement. L'atmosphère \emph{pousse} sur le système qui se
        contracte : celui-ci \rep{reçoit} du travail, d'où $W>0$.
\end{enumerate}
\end{correction}

\begin{correction}[retrouver le travail échangé]
\begin{enumerate}[label=\alph*),leftmargin=2em,topsep=2pt,itemsep=2pt]
  \item $\dU = Q + W \Rightarrow W = \dU - Q = 500 - 800 = \rep{\SI{-300}{\joule}}$.
  \item $W<0$ : le gaz \rep{cède} du travail ; il s'est donc \rep{dilaté}. (Il reçoit \SI{800}{\joule} de
        chaleur, en garde \SI{500}{\joule} en énergie interne et en restitue \SI{300}{\joule} sous forme de travail.)
\end{enumerate}
\end{correction}

\begin{correction}[placer les produits sur le diagramme]
On utilise $H_{\text{produits}} = H_{\text{réactifs}} + \dH$.
\begin{enumerate}[label=\alph*),leftmargin=2em,topsep=2pt,itemsep=2pt]
  \item $H_{\text{produits}} = 60 + 150 = \rep{\SI{210}{\kilo\joule}}$ : produits \emph{au-dessus},
        réaction \rep{endothermique}.
  \item $H_{\text{produits}} = 60 - 45 = \rep{\SI{15}{\kilo\joule}}$ : produits \emph{en dessous},
        réaction \rep{exothermique}.
  \item C'est le cas \rep{(a)} : $\dH>0$, la flèche va des réactifs (bas) vers les produits (haut), donc
        vers le \emph{haut}.
\end{enumerate}
\end{correction}

\begin{correction}[lire le « + Énergie » d'une équation]
\begin{enumerate}[label=\alph*),leftmargin=2em,topsep=2pt,itemsep=2pt]
  \item Le « $+\,$Énergie » côté produits signale que la réaction \emph{libère} de l'énergie : elle est
        exothermique, donc \rep{$\dH<0$}.
  \item Cette énergie provient de l'\rep{énergie chimique stockée dans les liaisons} des réactifs,
        restituée lors de l'oxydation du carbone et de l'hydrogène.
  \item Puisque $\dH = H_{\text{produits}} - H_{\text{réactifs}} < 0$, on a $H_{\text{produits}} <
        H_{\text{réactifs}}$ : les \rep{réactifs} ont l'enthalpie la plus \emph{élevée}.
\end{enumerate}
\end{correction}

\end{document}
